\chapter{CEA (Carcinoembryonic Antigen)}

\section*{What Is This Test?}
Carcinoembryonic antigen (CEA) is a 180--200 kDa oncofetal glycoprotein that is a member of the immunoglobulin superfamily of cell adhesion molecules. CEA is normally produced during fetal development (principally by the gastrointestinal tract, liver, and pancreas) and its levels fall to very low or undetectable levels after birth. In adults, CEA is present at low concentrations in normal colon, gastric, and bronchial epithelium. Elevated CEA levels are found in several malignancies, most notably colorectal adenocarcinoma (60--70\% of patients with advanced disease), and also in pancreatic, gastric, breast, lung, medullary thyroid, and ovarian cancers. CEA is the most established and widely used serum tumor marker for colorectal cancer, where its primary role is in monitoring for recurrence after curative surgical resection and monitoring treatment response in metastatic disease \cite{duffy2007tumour}. CEA is not sensitive or specific enough to be used as a screening test for colorectal cancer in asymptomatic individuals \cite{fletcher1986carcinoembryonic}. Benign conditions---particularly cigarette smoking, inflammatory bowel disease, pancreatitis, liver disease, and any condition causing hepatic inflammation or biliary obstruction---can elevate CEA levels.

\section*{Alternative Names}
CEA; Carcinoembryonic Antigen; CEA Blood Test; Colorectal Cancer Tumor Marker; Oncofetal Antigen

\section*{Principle}
CEA is measured by a two-site sandwich chemiluminescent immunoassay (CLIA) or electrochemiluminescent immunoassay (ECLIA) on automated immunoassay analyzers. Monoclonal antibodies directed against distinct epitopes on the CEA molecule are used as capture and detection antibodies. Patient serum is incubated with both antibodies. The capture antibody is typically biotinylated and binds to streptavidin-coated paramagnetic microparticles. The detection antibody is conjugated to a chemiluminescent label (ruthenium complex, acridinium ester, or horseradish peroxidase). After incubation and washing, a chemiluminescent or electrochemiluminescent signal is generated, and the light emission is proportional to the CEA concentration. Results are calibrated against the WHO 1st International Reference Preparation for CEA (IRP 73/601) and reported in ng/mL (or mcg/L, which is equivalent). CEA is highly stable in serum. Inter-method variability between platforms exists; serial CEA measurements should be performed using the same assay and ideally the same laboratory. Heterophile antibodies and HAMA (human anti-mouse antibodies) can interfere and cause false-positive elevation.

\section*{History}
CEA was discovered in 1965 by Phil Gold and Samuel Freedman at McGill University \cite{gold1965demonstration}. They extracted CEA from human colorectal adenocarcinoma tissue and from fetal gastrointestinal tissue (hence the name ``carcinoembryonic''). They demonstrated that CEA was present in fetal gut, liver, and pancreas during the first two trimesters but absent from normal adult tissues. In 1969, Thomson and colleagues developed the first radioimmunoassay (RIA) for CEA. Throughout the 1970s, CEA was extensively studied as a potential screening test for colorectal cancer, but its lack of sensitivity for early-stage disease and its non-specific elevation in benign conditions limited its screening utility. The American Society of Clinical Oncology (ASCO) published clinical practice guidelines for the use of CEA in colorectal cancer in 1996, 2000, and 2015, establishing its role in post-operative surveillance and monitoring \cite{locker2006asco}. CEA monitoring is now a standard component of the College of American Pathologists (CAP) and NCCN guidelines for colorectal cancer.

\section*{Why This Test Is Ordered}
\begin{itemize}
  \item Post-operative surveillance for colorectal cancer recurrence after curative-intent surgical resection --- CEA should be measured every 3--6 months for the first 3 years, then every 6 months for years 4--5 per NCCN guidelines \cite{benson2021colon}. A rising CEA after a period of normal or stable levels is the most sensitive indicator of recurrence, often predating clinical or radiographic evidence by 3--8 months
  \item Preoperative CEA measurement in newly diagnosed colorectal cancer for prognostication and as a baseline for post-operative monitoring --- elevated preoperative CEA is associated with larger tumor size, higher stage, and poorer prognosis
  \item Monitoring treatment response in metastatic colorectal cancer during chemotherapy --- declining CEA indicates treatment response; rising CEA indicates disease progression
  \item Monitoring of medullary thyroid cancer (along with calcitonin) and gastrointestinal neuroendocrine tumors
  \item Adjunct in the evaluation of solitary pulmonary nodules or pleural effusions --- elevated CEA in pleural fluid relative to serum (pleural fluid/serum CEA ratio >1.0) suggests malignant pleural effusion (primary lung adenocarcinoma or metastatic adenocarcinoma)
  \item Monitoring of pancreatic adenocarcinoma (supplementary to CA 19-9) and gastric adenocarcinoma
  \item Monitoring of breast cancer with documented CEA elevation that correlates with tumor burden (supplementary to CA 15-3/CA 27.29)
\end{itemize}

\section*{Primary vs Secondary Test}
CEA is a secondary test. It is NOT recommended for colorectal cancer screening in asymptomatic individuals. Its primary role is in monitoring for recurrence and treatment response in patients with established colorectal cancer.

\section*{How This Test Is Performed}
For most patients, a health care professional will take a blood sample from a vein in your arm using a small needle. After the needle is inserted, a small amount of blood will be collected into a test tube or vial. You may feel a little sting when the needle goes in or out. This usually takes less than five minutes.

\section*{What Preparation Is Needed}
No special preparation is required. Fasting is not necessary. CEA levels are higher in cigarette smokers (median 2--3 times higher than in non-smokers), and patients should be asked about smoking status when interpreting results. For patients being followed for colorectal cancer recurrence, CEA should ideally be measured at the same laboratory using the same assay to minimize inter-assay variability.

\section*{Contraindications and Precautions}
\begin{itemize}
  \item No absolute contraindications. The most important pre-analytical consideration is that CEA is non-specific and is elevated in a wide range of benign conditions: cigarette smoking (the most common cause of mild CEA elevation; CEA is typically 2--5 ng/mL in smokers, upper limit 7--10 ng/mL; smoking cessation reduces CEA to non-smoker baseline over 6--12 months), chronic obstructive pulmonary disease, inflammatory bowel disease (ulcerative colitis, Crohn disease with active colonic inflammation; CEA often 5--15 ng/mL), acute and chronic pancreatitis, cirrhosis and chronic hepatitis, obstructive jaundice (any cause), peptic ulcer disease, colonic polyps, diverticulitis, and benign breast disease. Hepatic metabolism and biliary excretion are the main routes of CEA clearance --- any process that impairs hepatic function (cirrhosis, metastatic liver disease, cholestasis) can significantly elevate CEA due to impaired clearance. Specimen should be collected in a serum separator tube (SST, gold-top) or red-top tube. Approximately 20--30\% of patients with colorectal adenocarcinoma (especially early-stage disease) have normal CEA levels, and a normal CEA does not exclude recurrence.
\end{itemize}
\section*{Risks}
Minimal risk. Slight pain or bruising at the venipuncture site. Rare: hematoma, infection, vasovagal syncope.

\section*{What Happens After the Test}
Results are typically available within 1 to 2 hours. An elevated preoperative CEA in a patient with newly diagnosed colorectal cancer should be documented as a baseline value. After curative-intent surgical resection, CEA should normalize within 4--6 weeks. Failure to normalize suggests residual disease (incomplete resection). A rising CEA during post-operative surveillance should prompt a thorough investigation for recurrence with colonoscopy (to detect anastomotic recurrence), CT of the chest, abdomen, and pelvis, and consideration of FDG-PET/CT when imaging is negative but CEA continues to rise.

\section*{Understanding Results}

\subsection*{Below Normal}
\begin{itemize}
  \item CEA <3 ng/mL in non-smokers (<5 ng/mL in smokers): Normal range. Does not exclude colorectal cancer, as approximately 20--30\% of patients with colorectal adenocarcinoma (particularly early-stage) have normal CEA levels
  \item Normal CEA after curative-intent surgical resection for colorectal cancer: Indicates complete resection (R0). CEA should normalize within 4--6 weeks after surgery
  \item Normalization of elevated CEA during chemotherapy for metastatic colorectal cancer: Indicates treatment response. A decrease of $\geq$50\% during therapy is associated with favorable prognosis
  \item CEA non-secretors: A subset of patients with colorectal cancer never express elevated CEA despite significant tumor burden. In these patients, CEA monitoring is not useful, and post-operative surveillance relies on imaging and colonoscopy
\end{itemize}

\subsection*{Above Normal}
\begin{itemize}
  \item CEA >5 ng/mL in non-smokers (>10 ng/mL in smokers): Elevated and warrants investigation if there is a history of colorectal cancer or if clinical findings suggest malignancy. CEA levels that do not return to normal after a documented elevation tend to rise progressively in the setting of untreated cancer
  \item CEA 5--15 ng/mL: Mild to moderate elevation. Common causes include cigarette smoking (the most common cause), inflammatory bowel disease, benign liver disease (cirrhosis, hepatitis, obstructive jaundice), benign colonic polyps, chronic pancreatitis, and early-stage colorectal cancer. Clinical correlation is essential
  \item CEA >15--20 ng/mL: More strongly associated with malignancy than benign disease. Elevation to this level in a patient with a history of colorectal cancer is highly suspicious for recurrence
  \item CEA >50 ng/mL: High probability of malignancy and metastatic disease. In the absence of a prior cancer diagnosis, warrants thorough investigation including colonoscopy and CT chest/abdomen/pelvis
  \item CEA >500--1,000 ng/mL: Almost always indicates extensive disseminated malignancy, most commonly metastatic colorectal adenocarcinoma with hepatic metastases. At these levels, the liver is almost always involved
  \item Rising CEA during post-operative surveillance: A confirmed rise in CEA from a previously normal or stable value is the single most sensitive indicator of colorectal cancer recurrence. A CEA elevation >2 ng/mL above the baseline or >5 ng/mL absolute value should trigger investigation for recurrence. CEA typically rises 3--8 months before recurrence is evident on CT or clinical examination. The rate of rise (CEA doubling time) is prognostic --- a CEA doubling time <30--60 days is associated with aggressive disease biology and poor prognosis
  \item CEA elevation in non-colorectal malignancies: Pancreatic adenocarcinoma (40--60\%), gastric adenocarcinoma (40--50\%), medullary thyroid cancer (CEA is a mandatory tumor marker together with calcitonin), breast cancer (30--50\%, supplementary to CA 15-3), lung adenocarcinoma (50--70\%), ovarian mucinous carcinoma, and malignant pleural effusion (pleural fluid/serum CEA ratio >1.0)
  \item False-positive CEA elevation: Smoking (the most common, 2--5 ng/mL), inflammatory bowel disease (colitis), cirrhosis and chronic hepatitis, acute pancreatitis, biliary obstruction, and recent surgery or trauma
\end{itemize}

\section*{Normal Results}
\begin{itemize}
  \item \textbf{CEA (non-smokers):} <3.0 ng/mL (for non-smoker adults). Some laboratories use <2.5 ng/mL for non-smokers and <5.0 ng/mL for smokers
  \item \textbf{CEA (smokers):} <5.0 ng/mL (can range up to 7--10 ng/mL in heavy smokers). Smoking cessation reduces CEA to non-smoker levels within 6--12 months
  \item \textbf{CEA (age >65):} Values tend to be slightly higher in the elderly
  \item \textbf{CEA in healthy individuals:} <2.5 ng/mL in >97\% of healthy non-smokers
  \item \textbf{CEA trend:} More important than a single value. A progressive rise over 2--3 serial measurements is more concerning than a single borderline elevation, which may represent benign or transient causes
  \item Serial CEA measurements should be performed using the same assay and ideally the same laboratory, as inter-assay and inter-laboratory variability can be significant (CV 10--15\%)
\end{itemize}

\section*{Follow-up Tests}
\begin{itemize}
  \item \textbf{Colonoscopy:} The definitive test for detecting colorectal cancer recurrence at the anastomotic site or metachronous colorectal cancers. Should be performed 1 year after surgical resection, then at 3 years, then every 5 years if normal (per NCCN and US Multi-Society Task Force guidelines)
  \item \textbf{CT chest/abdomen/pelvis with intravenous and oral contrast:} For detection of locoregional recurrence and distant metastases (liver, lung, peritoneum, lymph nodes) in patients with rising CEA. Recommended every 6--12 months for the first 3--5 years after curative-intent colorectal cancer surgery
  \item \textbf{FDG-PET/CT (fluorodeoxyglucose positron emission tomography):} For patients with unequivocally rising CEA on at least two occasions despite negative conventional imaging (CT and colonoscopy). PET-CT detects metabolically active recurrent disease that may be occult on CT, particularly peritoneal and nodal metastases. Highly sensitive for detecting sites of recurrence
  \item \textbf{Contrast-enhanced MRI of the liver:} More sensitive than CT for detecting and characterizing hepatic metastases. Hepatic resection or ablation of isolated liver metastases from colorectal cancer can be curative in selected patients
  \item \textbf{Diagnostic laparoscopy:} For patients with rising CEA and negative noninvasive imaging, peritoneal exploration can detect peritoneal carcinomatosis not visible on CT
  \item \textbf{Calcitonin:} If medullary thyroid cancer is suspected alongside CEA elevation. Both calcitonin and CEA are tumor markers for medullary thyroid cancer
  \item \textbf{CA 15-3 or CA 27.29:} If a primary breast cancer is suspected and CEA is the sole elevated marker. CA 15-3 is the primary serum marker for breast cancer monitoring
  \item \textbf{CA 19-9:} If pancreatic adenocarcinoma is suspected in the setting of elevated CEA, as dual elevation of CEA and CA 19-9 occurs in pancreatic cancer
\end{itemize}

\section*{References}
\bibliographystyle{unsrtnat}
\bibliography{references}

\clearpage
\section*{Regulatory Status and Authorization Date}
This chapter describes a test category, not a specific commercial product. FDA, CDSCO, EU IVDR, and MHRA approval or registration dates therefore cannot be assigned without the assay manufacturer, product name, instrument, and intended use. For laboratory-developed tests, an individual agency approval date may not exist. See the Preface for the interpretation of regulatory status.

\clearpage
